
Tools for assessing machine learning reproducibility---statistical evaluators, LLM-based pipeline screeners, end-to-end verifiers---share a fundamental limitation: they operate after experiments have already run. A researcher selecting a benchmark dataset receives no signal about whether their chosen dataset will yield reproducible results until they have invested substantial computational resources and development time. This reactive paradigm persists despite widespread recognition that reproducibility failures trace to upstream causes, particularly ambiguous dataset specifications that permit divergent preprocessing implementations.

This work investigates whether dataset metadata completeness predicts reproducibility variance \textit{before} any experiment executes. Using matched experimental configurations from synthetic data modeled on OpenML benchmark datasets (identical flows and hyperparameters), results indicate that top-quartile metadata completeness is associated with a 42.1\% reduction in interquartile range (IQR) of performance outcomes compared to bottom-quartile datasets (95\% CI: 39.1--51.7\%).

\subsubsection{The Problem of Post-Hoc Assessment}

The machine learning reproducibility crisis is documented in prior work. Kapoor and Narayanan (2022) identified eight types of data leakage affecting 329 papers across 17 scientific fields. The research community responded with assessment infrastructure: rliable \citep{agarwal2021deep} provides statistical tools for reliable benchmark evaluation; Reproscreener \citep{bhaskar2024reproscreener} uses language models to assess pipeline reproducibility; paper-replay enables end-to-end verification with cryptographic attestation. These tools share a common characteristic: they verify reproducibility \textit{post-hoc}, after experiments complete.

Yet reproducibility variance may originate upstream---in the ambiguities that dataset documentation leaves unresolved. When metadata omits preprocessing specifications, missing value handling, or train/test split definitions, researchers facing ambiguous datasets may make divergent implementation choices. This implementation heterogeneity propagates through the experimental pipeline, potentially manifesting as variance in reported results.

\subsubsection{Dataset Metadata as an Upstream Predictor}

This work examines a conceptual reframing: reproducibility as a continuous property of datasets, potentially predictable from metadata characteristics before any experiment runs. The proposed mechanism is epistemic entropy reduction---complete documentation constrains the degrees of freedom available to implementing researchers, reducing the space of valid preprocessing pipelines and thereby reducing outcome variance.

Mediation analysis in this study indicates that preprocessing entropy accounts for 64.7\% of the metadata-to-variance association (Sobel Z = 16.02, p < 0.0001). Datasets with more complete metadata exhibit lower Shannon entropy in preprocessing component distributions (imputation methods, scaling approaches, encoding strategies). This lower entropy corresponds to reduced performance variance across matched runs.

\subsubsection{Contributions}

This work makes three contributions:

First, it provides evidence that metadata completeness is associated with reproducibility variance. A 5-field metadata checklist---train/test split specification, preprocessing enumeration, missing value handling documentation, feature semantics provision, and versioning presence---accounts for variance in IQR after controlling for intrinsic dataset stability, popularity, algorithm family, and infrastructure factors.

Second, preprocessing entropy is identified as a candidate mediating mechanism. The 64.7\% mediation proportion suggests that documentation completeness may operate primarily by constraining preprocessing choices rather than through alternative mechanisms.

Third, robustness checks address potential confounds. The association persists in early-run subsamples (90.7\% preservation, addressing reverse causality concerns), within single algorithm families (p < 0.001 for RandomForest-only analysis, addressing algorithm-mix confounding), and against permutation controls (observed effect exceeds 95th percentile of null distribution).

\subsubsection{Limitations}

All experiments were conducted on synthetic data due to OpenML API availability constraints (504 Gateway Timeout during data collection). While the synthetic data follows expected distribution patterns, validation with real-world data is necessary before drawing conclusions about actual OpenML datasets. The design is observational---metadata completeness cannot be randomized---so causal claims are not supported. Results are specific to the synthetic data modeled on OpenML tabular datasets; generalization to other platforms requires separate validation.

